\chapter*{Introduction generale}
\addcontentsline{toc}{chapter}{Introduction generale}

Dans un contexte touristique de plus en plus numerise, les voyageurs attendent des reponses rapides, personnalisees et multilingues. Les hotels doivent ainsi offrir des experiences digitales modernes capables d'ameliorer la relation client avant, pendant et apres le sejour.

Notre projet de fin d'etudes s'inscrit dans cette dynamique. Il vise a concevoir une plateforme web intelligente permettant aux hotels tunisiens de centraliser leurs informations, d'assister les visiteurs via un chatbot base sur l'IA generative et de suivre les interactions utilisateurs a travers un module analytics.

La solution proposee repose sur une architecture moderne : Next.js pour l'application web, PostgreSQL pour la persistance des donnees, Redis pour le cache et les sessions, ainsi qu'une integration LLM pour la generation de reponses contextualisees via une approche RAG.

Afin d'atteindre ces objectifs, ce rapport est organise comme suit :
\begin{itemize}
  \item Le premier chapitre presente le contexte general du projet.
  \item Le deuxieme chapitre detaille l'analyse des besoins et l'architecture.
  \item Le troisieme chapitre expose la conception et la modelisation.
  \item Le quatrieme chapitre decrit l'implementation du chatbot intelligent.
  \item Le cinquieme chapitre traite le module analytics et l'espace administrateur.
  \item Le sixieme chapitre couvre la securite, les tests et le deploiement.
\end{itemize}
